\chapter{Results}
In what follows, we present three sets of results. First we studied the baseline model estimating the effect of the trade exposure index in conflict taking into account the economic heterogeneity. Second we show the model taking into account the geographical, sectorial and tribal heterogeneity. Third, we investigate the effect of the index across region with different water availability. \\

\section{Baseline model}
Table\ref{tab:ports_main} reports the estimates of the main coefficient from Equation \eqref{eq1} and \eqref{eq2}.
\begin{table}[htbp]
	\centering
	\caption{Correlation between port shocks and conflict: commercial vs all ports}
	\label{tab:ports_main}
	\begingroup
	\def\sym#1{\ifmmode^{#1}\else\(^{#1}\)\fi}
	\begin{adjustbox}{max width=\textwidth}
		\begin{tabular}{lcccc}
			\toprule
			& (1) Commercial ports & (2) Commercial ports & (3) All ports & (4) All ports \\
			\midrule
			Trade Exposure Index        & -1.01e-07    &             &        -9.22e-09     &             \\
			& (6.30e-08)   &             &   (2.1129e-08)       &             \\
			\addlinespace[0.25em]
			*High Nightlight Intensity       &              & -1.12e-07   &             &   -3.12e-08             \\
			(1 = Above the median)&              & (7.80e-08)  &             &           (2.85e-08)  \\
			*Low Nightlight Intensity        &              & -5.61e-08\sym{***} &       & 7.30e-09           \\
			(1 = Below the median)			&              & (5.43e-09)  &             &    (7.51e-09)         \\
			\addlinespace[0.25em]
			
			\addlinespace[0.25em]
			Nightlight intensity     & -0.879\sym{***} & -0.878\sym{***} & -0.902\sym{***} & -0.901\sym{***} \\
			& (0.249)     & (0.248)     & (0.259)     & (0.260)     \\
			\midrule
			Observations         & 1,611       & 1,611       & 1,611       & 1,611       \\
			Time fixed effect     & YES       & YES       & YES       & YES      \\
			District fixed effect           & YES       & YES       & YES       & YES       \\
			R² (Adj.)          & 0.662       & 0.662       & 0.661       & 0.662       \\
			\hline
		\end{tabular}
		\end{adjustbox}
		\vspace{3pt}
		\\[2pt]
		\begin{minipage}{1\textwidth}
			\footnotesize Fixed effects: ADM2 and year. SE clustered at ADM2 level. \\
			High/Low Nightlight Intensity are ADM2 groups defined by the sample median of nightlights (harmonized DMSP-OLS/VIIRS). Entries in group rows report the simple slope of the Trade Exposure Index within the indicated group; group dummies themselves are absorbed by ADM2 FE. 358 observations dropped due to missing RHS variables. Signif.: *** p < 0.01, ** p < 0.05, * p < 0.10.
		\end{minipage}
	\endgroup
\end{table}

It shows the results of the baseline specification, comparing the commercial ports with the broader set of Nigerian ports. 
Column (1) and (2) focus exclusively on the commercial ports, while Columns (3)–(4)exclude them as a falsification test. \\ 
Starting from Column (1), results indicates that the aggregate coefficient on the exposure index is small and not not statistically significant. However, when the sample is split to study the heterogeneity between high and low nightlight areas, as in Column (2), a strong heterogeneity emerges. This division was made to distinguish ADM2 units by local economic development, proxied by night time lights (as detailed in the Data section): where the areas with higher income are the ones with values above the 75 percentile \parencite{amodio2024trade, henderson2012measuring}. \\
In the areas with reatevely low levels of nightlights, the effect of the commercial shock is negative and statistically significant, suggesting that a greater exposure maritme trade reduces the likelihood of consumption based conflicts. Quantitatively, the coefficient implies that the trade index reduce conflict by 6 percent. \\ 
In contrast the coefficient is close to zero in the more developed area. This model is coherent with the intuition that the external trade shocks have stronger effect where the local economies are less diversified and more vulnerable to changes in global market integration \parencite{berman2015external}. Thus distance from commercial ports is particularly important in the less developed regions, amplifying the regulatory effect of the trade shocks. This finding is especially relevant given that less developed areas are also more likely to have conflict \parencite{collier2004greed, bruckner2010international} \\

Columns (3)–(4) serve as a falsification test. If the relationship was determined solely by proximity with any port, not only the commercial one, we would expect similar effects when Nigerian trade ports are excluded.
Columns 3 and 4 extend the analysis, excluding the main trade ports. The results show weaker and not statistically significant coefficients, both for the aggregate and the splitted samples.
The absence of an effect reinforces the interpretation that the mechanism is specificaly linked to international trade flows, rather than mere geographic proximity. \\

Overall, the results suggest that the link between trade and conflict is neither mechanic nor omnipresent, it rather depends commercial relevance of ports and the level of local economic development. The coefficients support the idea that international integration can provide stabilising benefit in the peripherical regions, and so the less developed ones economically speaking, while its marginal impact in more developed areas remains limited.\\
\section{Heterogeneity}

\begin{table}[htbp]
	\caption{Regional and sectoral heterogeneity (combined)}
	\label{tab:hetero_all}
	
	\scriptsize
	\setlength{\tabcolsep}{4pt}
	\renewcommand{\arraystretch}{1.05}
	\textbf{Panel A. Index effects} \\
	% --- PANEL A ---
	\resizebox{1\textwidth}{!}{
		\begin{tabular}{lc|c|c|cc}
			\hline
			& East vs West
			& Chemicals vs Vegetables
			& Export vs Import
			& \multicolumn{2}{c}{Tribal} \\
			\hline
			
			East (slope) & -1.63e-07 *  & & &&\\
			& (8.24e-08)   &   &                  &&  \\
			West (slope) &  -7.45e-09 & &&&\\
			&    (4.76e-08)   & &          &     &  \\
			Vegetables (slope)  &  & -1.63e-07 * & && \\
			&  & (8.24e-08) & && \\
			Chemicals (slope) &  & -7.45 e-09 && &\\
			& &   (4.76e-07) & && \\
			Export (slope)   & & & 2.13e-09 &&\\
			 & & & (3.29e-08)& &\\
			Import (slope)  & & &  -2.55e-07 *** &&\\
			 & & &  (4.87e-08) && \\
			Low Nightlight Intensity         &&& & 5.55e-05           & -5.7e-08*** \\
			(1 = Below the median)  &&&&(4.41e-05)      & (5.43e-9) \\
			High Nightlight Intensity       &&&   & -3.48e-07    &-1.63e-08 \\
			(1 = Above the median) &&&& (2.05e-07)   & (-9.34e-09) \\
			Dominance x Low NTL   &&& & 5-53 e-05        & \\
			 &&&&(4.13e-05 )       & \\
			Dominance x High NTL    &                 &&& 5.16e-07 &\\
			 &                 &&&    (2.77e-09) &\\
			Fractionalization x Low NTL &&&&& 3.88e-05\\
			&&&&&(2.87e-05)\\
			Fractionalization x High NTL &&&&& 3.42e-07\\
			&&&&&(2.45e-07)\\
			\addlinespace[2pt]
			\textit{Controls} \\
			Nightights intensity & -0.884 *** & -0.884 *** & -0.884 *** & -0.877*** & -0.877*** \\
			 & (0.2493) & (0.2493) & (0.2470) &(0.248) & (0.248) \\
			\hline
			Observations 1,611 & 1,611 &1,611 & 1,611 & 1,611 & 1,611 \\
			Time fixed effect &YES& YES & YES & YES & YES\\
			District fixed effect&YES& YES & YES & YES & YES\\
			Adj. R$^2$    &0.665& 0.665 &0.670 & 0.662 & 0.662\\
			\hline
		\end{tabular}
	} % fine resize Panel A
	
	\vspace{4pt}
	\textbf{Panel B. Average nightlights (mean / median / sd)}\\[2pt]
	% --- PANEL B (stessa larghezza del Panel A) ---
	\tiny
	\resizebox{0.8\textwidth}{!}{%
		\begin{tabular}{lcccccccc}
			\toprule
			Group && Mean NL && Median NL && SD NL && N \\
			\midrule
			East  && 9.78 && 6.57 && 10.4 && 979 \\
			West && 16.8 && 8.91 && 18.9 && 990 \\
			\midrule
			Chemicals && 16.5 && 8.41 && 19.2 && 957 \\
			Vegetables && 10.3 && 6.76 && 10.6 && 1,012 \\
			\midrule
			Export  && 16.4 && 9.17 && 17.9 && 1,199 \\
			Import  && 8.48 && 5.83 && 9.52 && 770 \\
			\bottomrule
		\end{tabular}
	} % fine resize Panel B
	\vspace{3pt}
	\\[2pt]
	\begin{minipage}{0.95\textwidth}
		\footnotesize \textbf{Panel A} Fixed effects: ADM2 and year. SE: clustered at ADM2 level. Rows labeled “(slope)” report the slope of the Trade Exposure Index within the stated group, estimated via specifications that allow group-specific slopes. High/Low Nightlight Intensity are ADM2 groups defined by the sample median of nightlights (harmonized DMSP-OLS/VIIRS). Tribal indices are time-invariant; their levels are absorbed by ADM2 FE, so only interaction driven slopes are identified. Signif.: *** p < 0.01, ** p < 0.05, * p < 0.10. \\
		\textbf{Panel B} Nightlights (NTL) are mean annual radiance from harmonized DMSP-OLS/VIIRS. Statistics are computed at the ADM2 level over 2009–2019 (Mean/Median/SD across years within ADM2; N = number of ADM2 in each group).\\
	\end{minipage}
\end{table}


After establishing the baseline relationship, we explore whether the estimated effects vary across Nigeria’s main macro regions and according to local social fragmentation. Table \ref{tab:hetero_all} combines these dimensions of heterogeneity: Panel A shows how the trade conflict relationship differs across regions, sectors and tribal fominance and fractorization, whereas Panel B reports the corresponding nightlight averages, illustrating the underlying economic disparities between groups.\\

\subsection{\textit{Geography}} 
Columns (1) investigates geographic heterogeneity: the index is interacted with East/West dummies defined by the median longitude of coastal ADM2 centroids (details in Appendix~\ref{app:descriptives}). The estimated effect is more negative and statistically stronger in the eastern belt—areas that are poorer and more connected to maritime activity, consistent with a vulnerability to trade mechanism.
A detailed discussion of the spatial correlation between the macro regional partition, together with corresponding maps, is reported in Appendix~\ref{app:descriptives}.\\

These patterns mirror the spatial distribution of economic development (Panel B). Nightlight data shows that 
western ADM2s exhibit higher and more dispersed nightlight values, largely driven by the Lagos metropolitan area, while eastern units are uniformly dimmer, reflecting a more rural and less industrialized economic structure. In this case, the gap between richer and poorer areas (as proxied by nightlights) is particularly pronounced.
Detailed maps, descriptive statistics of nightlights by region, and additional robustness checks are reported in Appendix~\ref{app:descriptives}\\

\subsection{\textit{Sector}}
Columns (2)-(3) reports the sector heterogeneity: the index is interacted with two different sets of dummies: (i) the predominant type of industry traded through each port, either chemical or vegetable products (the two main sectors after minerals, based on data from \textcite{imfportwatch2025}, available since 2019 and used here as a proxy); and (ii) the predominant trade flow orientation, that is, whether the port handles mainly imports or exports (also based on data from IMF Portwatch). The estimated effect is negative and statistically significant for vegetables and imports. For presentation purposes, we refer to the coefficients as $\text{index}\_{veg}$ and $\text{index}\_{chem}$ (or $\text{index}\_{exp}$ and $\text{index}\_{imp}$), where the former represents the baseline effect and the latter includes the interaction term.
The stronger and statistically significant effects observed for vegetable and import oriented ports suggest that not only proximity and structural economic conditions matter, but also the nature and direction of trade flows. The results are consistent with previous findings that trade in basic consumption goods, such as agricultural products, can mitigate conflict by stabilizing local prices and incomes (\parencite{mcguirk2020economic}). The effect for vegetable oriented ports is also  consistent with what is shown by \textcite{amodio2024trade}: regions with lower nightlight intensity, typically linked to agricultural activity, are more vulnerable to economic shocks, suggesting that the welfare channel of trade is particularly relevant in such contexts. Similarly, import flows—particularly of essential goods—may enhance welfare and reduce consumption related tensions \parencite{bruckner2010international, dube2013commodity}. Overall, these patterns reinforce the idea that trade integration contributes to stability primarily through welfare improving, rather than resource based, channels.\\
Results are aligned with the underlying economic structure summarized in Panel B, where vegetable and import oriented ports correspond to areas exhibiting more dispersed nightlight distributions, a pattern typically associated with more rural and less industrialized economies. Although the import/export distinction also reflects structural differences in local economies, import oriented ports show a systematically stronger effect than all other categories. Overall, these patterns reinforce the idea that trade integration contributes to stability primarily through welfare improving, rather than resource based, channels. \\

\subsection{\textit{Tribes}} 
Columns (5)–(6) report the results for the heterogeneity by ethnic composition.\\
Information on tribal distribution is used to proxy local ethnic diversity.
Two indicators were constructed: tribal dominance and tribal fractionalization.\\
Tribal dominance identifies areas where a single ethnic group constitutes more than 50 percent of the local population \parencite{esteban2012ethnicity}, while tribal fractionalization measures the probability that two randomly selected individuals belong to different ethnic groups \parencite{collier2004greed, collier1998economic}.\\
Empirical studies suggest that these two dimensions have opposite effects on conflict: \textcite{esteban2012ethnicity} find that higher fractionalization tends to reduce conflict, whereas \textcite{collier2004greed} show that dominance may fuel conflict through capture of local rents and exclusion of minorities.\\
In our estimations, both coefficients are positive but not statistically significant, indicating no robust relationship between trade exposure and conflict across different levels of ethnic composition.\\
However, when accounting for ethnic fractionalization, the coefficient on the trade index in low nightlight areas remains negative and stastically significant, consistent with the baseline results in Table~\ref{tab:ports_main}.
\subsection{\textit{Water avaibility}}
\begin{table}[htbp]\centering
	\caption{Water Availability: Levels \& Heterogeneity}
	\label{tab:water}
	\begin{adjustbox}{max width=\textwidth}
		\begin{tabular}{lr}
			\toprule
			& Terciles \\
			\midrule
			Index × Water tercile: Low & 6.48e-08 \\
			(1 = below 33\%) & (8.76e-08)  \\
			Index × Water tercile: Medium & -2.597e-07\sym{**} \\
			(1 = between 33\% - 66\%)&  (8.75e-08) \\
			Index × Water tercile: High & 3.78e-08 \\
			(1 = above 66\%) &  (9.38e-08) \\
			Nightights intensity & -0.834\sym{***} \\
			& (0.239) \\
			\midrule
			Observations  & 1620 \\
			Time Fixed effect & YES\\
			District Fixed Effect & YES\\
			Adj. $R^2$  & 0.669 \\
			\bottomrule
		\end{tabular}
	\end{adjustbox}
	\begin{flushleft}\footnotesize
		Water availability is the ADM2 area‐weighted mean of a FAO raster, then split into terciles (cut at the 33rd/66th percentiles: Low/Medium/High). Because levels are absorbed by ADM2 fixed effects, we report Index slopes by water group (interactions). All models include year \& ADM2 FE with SE clustered at ADM2. 
		\sym{***} $p<0.01$, \sym{**} $p<0.05$, \sym{*} $p<0.10$.
	\end{flushleft}
\end{table}


We next investigate whether the pacifying effect of trade exposure varies with local resource endowments. Water availability, derived from the FAO–GAEZ database \parencite{FAO_GAEZ_2020_dataset}, serves as a proxy for agricultural potential and economic activity, given its close link to crop productivity and the prominence of vegetables trade in Nigerian ports.\\

The intuition is that in ADM2 units with lower water availability, typically more fragile and rural, the trade index should reduce conflict by stabilizing consumption and providing alternative income channels. This mechanism is consistent with evidence that economic integration can mitigate consumption related conflicts by improving access to markets and smoothing shocks \parencite{amodio2024trade}. \\

However, the empirical results in Table \ref{tab:water} reveal a more complex pattern. Contrary to expectations, the effect of trade exposure turns positive in areas with higher and low water availability, where agricultural and economic activities are already developed. This suggests that, rather than amplifying stability, trade may in resourch rich area vould reinforce existing inequalities or intensify competition over profitable markets \parencite{berman2017mine, michalopoulos2013pre}. Overall, the findings indicate that trade openness exerts a stabilizing influence mainly in less developed and resource constrained local economies, a result in line with broader evidence on the heterogeneous development impacts of market integration \parencite{dixcarneiro2023globalization, amodio2024trade}.